\chapter{Execution Beyond Equities}\label{ch:mx:execution-beyond-equities}

Buying the same amount of risk takes an order in an index future, a request for quote in bonds, a stream of last-look prices in currencies and a venue
that never closes in crypto. The objective is the same shortfall; the market's rules change what the algorithm can do. This chapter takes one risk
transfer, USD~50 million, through four simulated markets with execution adapters built on the market components of One Quant Books 1 to 3, and
measures the all-in cost in each and the part of it that the market's own rule explains.

The four markets are stylised: their spreads, volatilities and volumes (\cref{tab:mx:execution-beyond-equities:markets}) are assumptions of a
realistic order of magnitude, chosen to show the mechanisms, not measurements of particular instruments. Each order's impact follows the square-root
law of chapter~11, $Y\sigma\sqrt{Q/V}$ with $Y=0.7$, the daily volatility $\sigma$ and daily volume $V$.

\begin{table}[ht]
\centering
\small
\begin{tabular}{@{}lrrr@{}}
\toprule
market & half-spread (bp) & daily volatility (bp) & daily volume (USD) \\
\midrule
index future & 0.25 & 100 & 200 billion \\
EUR/USD & 0.08 & 50 & 300 billion \\
corporate bond & 5 & 40 & 20 million \\
bitcoin (one exchange) & 0.5 & 300 & 5 billion \\
\bottomrule
\end{tabular}
\caption{The simulated markets' parameters (assumptions). Data: \texttt{mx\_xexec.MARKETS}.}\label{tab:mx:execution-beyond-equities:markets}
\end{table}

\section{Futures: one book, implied liquidity, rolls}

An index future trades in one central book, with a tick that is often most of the spread and depth that dwarfs any single stock's. USD~50 million is
200 contracts at 5\,000 points and USD~50 a point. Its impact is small (1.11 basis points) and its spread smaller (0.25). What the future adds is the roll
(One Quant Book 1, chapter~21): a position held past expiry must be sold in the front month and bought in the next, and the calendar spread between
them has its own book (Book~1, chapter~19). Exchanges link the books with implied prices: a spread order can trade against the two outright books at
once (implied in), and spread orders make prices in the outrights (implied out); CME describes implied orders as removing the legging risk, and implied
quantity never matches implied quantity.

Rolling 200 contracts through this chapter's books (front 4\,999.75 bid for 150, back 5\,010.25 offered for 100, the spread book $-10.25$ bid for 60,
each level backed by the same size a tick worse): selling the front and buying the back as two orders costs 0.44 points a contract against the
spread's mid of $-10.25$, 0.88 basis points. Selling the spread in its own book for the first 60 contracts and legging the rest costs 0.23 points, 0.45
basis points: half. With implied prices the last 140 contracts trade at the same price as legging, but as one atomic order, with no risk that the
market moves between the legs (\texttt{firm\_xexec.roll}, on Book~1's \texttt{firm.match}).

\omcode{code/firm/xexec/firm_xexec.py}{76}{88}{Rolling a long futures position: legging through the two outright books, selling the spread in its own book, or using the book that implied prices show, then legging whatever is left.}\label{lst:mx:execution-beyond-equities:roll}

\section{FX: streams, last look and aggregation}

\begin{definition}[Liquidity aggregator, request for stream, firm liquidity]\label{def:mx:execution-beyond-equities:aggregator}
A \emph{liquidity aggregator}\index{liquidity aggregator} collects the streaming quotes of several liquidity providers into one book and routes each
request to the best price. A \emph{request for stream}\index{request for stream} asks providers to start streaming executable prices for a stated size
to one client. \emph{Firm liquidity}\index{firm liquidity} is liquidity whose quotes fill when hit, with no last-look window in which the provider may
reject.
\end{definition}

Most FX is traded bilaterally, off exchange, on providers' streams (One Quant Book 2, chapter~15). The FX Global Code's principle~17 asks providers
that use last look to be transparent about it and to use it as a risk control, not to gather information. The adapter sends USD~50 million as 20
requests of USD~2.5 million (\texttt{firm.acexec}'s straight-line schedule) to an aggregator of four last-look streams (half-spreads of 0.08 to 0.11
basis points, holds of 20 to 80 milliseconds, asymmetric thresholds of 0.05 to 0.08 basis points) and one firm stream at 0.15. A rejected request loses
the move that caused the rejection and goes on to the next stream at the new price (Book~2's \texttt{firm.lastlook}).

\omcode{code/firm/xexec/firm_xexec.py}{101}{113}{One request through the aggregator: the cheapest stream first; each rejection after its hold costs the client the move the provider declined to honour, and the request moves on.}\label{lst:mx:execution-beyond-equities:fx}

Over 2\,000 simulated orders, rejections numbered 7.7\% of the requests, and the aggregator's average cost was 0.087 basis points against the best
quote's 0.080: last look costs 0.007 basis points here, and the aggregator still beats the firm stream (0.15). Last look's cost is not in the spread
the client sees but in the moves it gives back on the requests it loses; Oomen (2017) analysed the practice in detail. With an impact of 0.45 basis
points, the all-in cost in FX is 0.54 basis points, the cheapest of the four.

\begin{dated}{2026-09}{How FX is traded}\label{dat:mx:execution-beyond-equities:bis}
The BIS Triennial Survey of April 2025 found electronic trading at 59\% of FX turnover, about the same share as in the previous survey, and noted that
customers using disclosed multi-dealer electronic trading could in theory choose among more than 15 platforms (BIS Quarterly Review, December 2025).
\end{dated}

\section{Bonds: requests for quote and portfolio lists}

\begin{definition}[Risk transfer price]\label{def:mx:execution-beyond-equities:rtp}
A \emph{risk transfer price}\index{risk transfer price} is a price at which a dealer takes a whole block at once as principal, charging for the cost and
risk of unwinding it, instead of working it as agent.
\end{definition}

A corporate bond trades a few times a day: USD~50 million is two and a half days of this bond's volume, and an order worked over a day would move the
price by the square-root law's 44.3 basis points on top of a 5-basis-point half-spread. The client instead asks dealers for a \omterm{def:mx:execution-beyond-equities:rtp}{risk transfer price} by
request for quote (Book~2, chapter~22). The adapter prices it with Book~2's \texttt{firm.rfq}: each dealer bids the value less a markup (here the
49.3 basis points the dealer expects to pay to unwind), plus a private term with a dispersion of 8 basis points (its inventory and axes); each dealer who
sees the request and loses costs the client 1.5 basis points of leakage. Asking one dealer costs 49.3; asking three, 45.5 (6.8 of competition less 3.0
of leakage); asking more adds leakage faster than competition, so three is best. Hendershott and Madhavan (2015) showed that such electronic
one-sided auctions are a viable source of liquidity even in inactively traded bonds.

\section{Crypto: round the clock, rate limits and on-chain legs}

Bitcoin trades on many exchanges and in automated market makers, every hour of every day. On one exchange with this chapter's parameters, USD~50
million costs 0.5 basis points of spread and 21.0 of impact. A constant-product pool (One Quant Book 3, chapter~20) with USD~1 billion on each side and
a 5-basis-point fee offers another leg: its price moves along $xy=k$ as the order takes coins out. The adapter sends $x$ to the pool and the rest to
the exchange and chooses $x$ to minimise the sum, the pool's price starting at the exchange's mid (\texttt{cex\_amm\_split}, on Book~3's integer
\texttt{firm.amm}): the best is USD~1.375 million to the pool, saving 0.36 basis points. The pool is cheap for the first dollars and expensive for
the rest: its marginal cost rises twice as fast as its average.

A venue's rate limits (Book~3, chapter~15) bound what the algorithm can send: a large spot exchange's published limits allow 100 orders in ten
seconds. Twelve \omterm{def:mx:the-almgren-chriss-framework:parent}{child orders} in an hour are nowhere near it; 120 in one second are refused (\texttt{orders\_allowed} on Book~3's
\texttt{firm.ratelimit}). Rate limits bind on strategies that quote and cancel, not on schedules like this one; the round-the-clock market does not
remove the day's volume pattern, it moves the thin hours to the night.

\section{One order, four markets}

\Cref{tab:mx:execution-beyond-equities:allin} puts the four executions side by side.

\begin{table}[ht]
\centering
\small
\begin{tabular}{@{}lrrrrl@{}}
\toprule
market & spread & impact & market's rule & all-in & the rule \\
\midrule
index future & 0.25 & 1.11 & 0.45 & 1.81 & roll through the spread book (legging: 0.88) \\
EUR/USD & 0.08 & 0.45 & 0.007 & 0.54 & last-look rejections \\
corporate bond & --- & 49.27 & $-3.77$ & 45.50 & RFQ to three dealers \\
bitcoin & 0.50 & 21.00 & $-0.36$ & 21.14 & a pool leg \\
\bottomrule
\end{tabular}
\caption{The all-in cost (basis points) of buying USD~50 million of risk in the four simulated markets, and the part of it that each market's rule
explains. For the bond, the dealer's markup takes the place of spread and impact. Data: \texttt{mx\_xexec.table}.}\label{tab:mx:execution-beyond-equities:allin}
\end{table}

The rule explains a quarter of the future's cost (the roll, 25\%), 1.3\% of the currency's, and it saves 8.3\% of the bond's (competition net of
leakage) and 1.7\% of bitcoin's. The costs themselves span two orders of magnitude, and the reason is the market's depth relative to the order, not its
rules: the same USD~50 million is a rounding error in the index future and two days' volume in the bond. An \omterm{def:mx:benchmark-algorithms:is}{implementation-shortfall algorithm}
(chapter~16) written against \texttt{firm.acexec}'s scheduler can drive all four through adapters; what changes is what a ``\omterm{def:mx:the-almgren-chriss-framework:parent}{child order}'' is: a
\omterm{def:mx:the-limit-order-book:orders}{limit order}, a request to a stream, an auction among dealers, a swap against a pool.

\section{Tutorial: one order, four markets}\label{tut:mx:execution-beyond-equities:tut}

\begin{tutorial}
\textbf{Goal.} Execute the same risk transfer through four execution adapters and split each all-in cost into spread, impact and the market's rule.
\textbf{End state:} \cref{tab:mx:execution-beyond-equities:allin} and the numbers of sections 1 to 4.
\begin{tutsteps}
\item \textbf{Futures.} \texttt{firm\_xexec.roll(front, back, spread, lots, tick)} on \texttt{firm.match}'s \texttt{Quote}; \texttt{mx\_xexec.futures()}.
\item \textbf{FX.} \texttt{LP}, \texttt{fx\_child}, \texttt{children(notional, scheduler, n)}; \texttt{fx()}.
\item \textbf{Bonds.} \texttt{rfq(n, markup, sigma, leak)}; \texttt{bonds()}.
\item \textbf{Crypto.} \texttt{amm\_cost\_bp}, \texttt{cex\_amm\_split}, \texttt{orders\_allowed}; \texttt{crypto()}; \texttt{table()}.
\end{tutsteps}
\textbf{What to change next.} Replace the bond's markup by a model of the dealer's own unwinding; give the FX streams symmetric last look; route the
bitcoin order across two exchanges and the pool.
\end{tutorial}

\section{Build: execution adapters}\label{bld:mx:execution-beyond-equities:xexec}

\begin{build}
\bfield{Purpose} The market-specific layer under chapter~28's \omterm{def:mx:benchmark-algorithms:algo}{execution algorithm}: one scheduler, several ways to trade a slice.
\bfield{Interface} \texttt{Market(name, half\_spread\_bp, sigma\_day\_bp, adv, y)}, \texttt{sqrt\_impact\_bp}, \texttt{children(notional,
scheduler, n)}, \texttt{roll(front, back, spread, lots, tick)}, \texttt{LP(name, half\_bp, hold\_ms, threshold\_bp, policy)},
\texttt{fx\_child(lps, sigma, rng)}, \texttt{rfq(n, markup, sigma, leak)}, \texttt{amm\_cost\_bp}, \texttt{cex\_amm\_split}, \texttt{orders\_allowed(governor,
horizon\_ms, n)}.
\bfield{Rules} Costs in basis points of notional against the starting mid; the adapters compose Books 1 to 3's components rather than copy them.
\bfield{Acceptance tests} \texttt{code/firm/xexec/tests/}: the roll on the chapter's books and implied-in at the legging price; a firm stream's cost is
its half-spread and an always-rejecting stream passes the request on; the last-look cost against Book~2's formula; the RFQ split adds up; a small
pool trade pays the fee; the split saves; the rate limit's boundary.
\bfield{Stretch} Futures in \texttt{firm.exchsim} with implied matching; symmetric last look; several pools and exchanges.
\end{build}

\begin{omsources}
\item Global Foreign Exchange Committee, \emph{FX Global Code}, principle~17 (last look).
\item R.~Oomen, ``Last look'', \emph{Quantitative Finance} 17(7), 2017.
\item T.~Hendershott and A.~Madhavan, ``Click or call? Auction versus search in the over-the-counter market'', \emph{Journal of Finance} 70(1),
2015.
\item CME Group, Globex matching algorithms and implied orders (client systems wiki).
\item BIS, ``The FX trade execution landscape through the prism of the 2025 BIS Triennial Survey'', \emph{BIS Quarterly Review}, December 2025.
\end{omsources}

\section{Exercises}

\begin{exercise}[$\star$]\label{exo:mx:execution-beyond-equities:1}
Compute the impact of USD~50 million in the index future and in the bond by the square-root law with the table's parameters.
\end{exercise}

\begin{exercise}[$\star$]\label{exo:mx:execution-beyond-equities:2}
Verify the roll's legging cost of 0.44 points a contract from the books.
\end{exercise}

\begin{exercise}[$\star$]\label{exo:mx:execution-beyond-equities:3}
What is the bond's cost by request for quote to one, two and four dealers?
\end{exercise}

\begin{exercise}[$\star\star$]\label{exo:mx:execution-beyond-equities:4}
Why does implied liquidity not improve the roll's price here, and what does it improve?
\end{exercise}

\begin{exercise}[$\star\star$]\label{exo:mx:execution-beyond-equities:5}
Why can an aggregator of last-look streams be cheaper than a firm stream, and when would it not be?
\end{exercise}

\begin{exercise}[$\star\star$]\label{exo:mx:execution-beyond-equities:6}
Why does the pool take only USD~1.4 million of the bitcoin order?
\end{exercise}

\begin{exercise}[$\star\star\star$]\label{exo:mx:execution-beyond-equities:7}
\emph{Coding.} With leakage of 3 basis points per losing dealer instead of 1.5, how many dealers should the client ask, and what does the bond cost?
\end{exercise}

\begin{exercise}[$\star\star\star$]\label{exo:mx:execution-beyond-equities:8}
\emph{Find the flaw.} ``Our FX algorithm paid 0.08 basis points on average: exactly the best quote in the aggregator, so last look cost us nothing.''
\end{exercise}

\section{Problem: One Order, Four Markets}

\begin{problem}[{Weekend problem --- one order, four markets}]\label{pb:mx:execution-beyond-equities:1}
A multi-asset fund must add USD~50 million of risk and can do it in an index future, a currency, a corporate bond or bitcoin. The desk asks what each
costs and why.

\medskip\noindent\textbf{Part I --- Futures.}
\begin{enumerate}
\item Why is a future's impact so small for this order?
\item Describe implied-in and implied-out prices.
\item Give the roll's cost by legging and through the spread book.
\item What does the implied book add?
\end{enumerate}

\medskip\noindent\textbf{Part II --- FX.}
\begin{enumerate}[resume]
\item Define an aggregator, a \omterm{def:mx:execution-beyond-equities:aggregator}{request for stream} and \omterm{def:mx:execution-beyond-equities:aggregator}{firm liquidity}.
\item What does principle 17 of the FX Global Code ask of last look?
\item How does the adapter model a rejection, and what does it cost?
\item Give the rejection rate and the last-look cost.
\end{enumerate}

\medskip\noindent\textbf{Part III --- Bonds and crypto.}
\begin{enumerate}[resume]
\item Define a \omterm{def:mx:execution-beyond-equities:rtp}{risk transfer price}.
\item Split the three-dealer RFQ cost into markup, competition and leakage.
\item Why is three dealers the best number here?
\item How is the bitcoin order split between the exchange and the pool, and what does it save?
\item When do an exchange's rate limits bind?
\end{enumerate}

\medskip\noindent\textbf{Part IV --- The comparison.}
\begin{enumerate}[resume]
\item \emph{State the named result}: the all-in cost of the same risk transfer in futures, FX, bonds and crypto in the simulator, and the part of it
each market's rule explains.
\item Why do the costs span two orders of magnitude?
\item What is a \omterm{def:mx:the-almgren-chriss-framework:parent}{child order} in each market?
\item How would you build one algorithm that trades all four?
\item Which of the chapter's assumptions would you check first against real data?
\item What does the dated box say about how FX is traded?
\item In one sentence: what changes when an \omterm{def:mx:benchmark-algorithms:algo}{execution algorithm} leaves the equity book?
\end{enumerate}
\end{problem}

\section{Interview questions}

\begin{interviewq}[$\star$ \iqroles{trader}]\label{iq:mx:execution-beyond-equities:1}
You must roll a large futures position. How do you do it, and what can go wrong?
\end{interviewq}

\begin{interviewq}[$\star\star$ \iqroles{trader}]\label{iq:mx:execution-beyond-equities:2}
A liquidity provider rejects 10\% of your FX requests. How do you work out what that costs you?
\end{interviewq}

\begin{interviewq}[$\star\star$ \iqroles{researcher}]\label{iq:mx:execution-beyond-equities:3}
How many dealers should a client ask for a price on a corporate bond, and what does the answer depend on?
\end{interviewq}

\begin{interviewq}[$\star\star$ \iqroles{developer}]\label{iq:mx:execution-beyond-equities:4}
Design an execution adapter interface that lets one scheduler trade equities, FX streams and RFQs.
\end{interviewq}

\begin{interviewq}[$\star\star$ \iqroles{researcher}]\label{iq:mx:execution-beyond-equities:5}
How would you split a large crypto order between a centralised exchange and an AMM pool?
\end{interviewq}

\begin{interviewq}[$\star\star\star$ \iqroles{bank}]\label{iq:mx:execution-beyond-equities:6}
A client asks your bank for a \omterm{def:mx:execution-beyond-equities:rtp}{risk transfer price} on USD~50 million of a bond that trades USD~20 million a day. How do you price it?
\end{interviewq}
